\booksection{Appendice: Analisi dell'espansione memoria su Android}

\chapterintro{Quando compri un cellulare, vedi spesso promozioni come “12GB+8GB RAM” — sembra quasi di avere 8GB di RAM in più dal nulla. Ma la rende davvero più veloce? E potrebbe persino danneggiare il cellulare? In questa sezione spieghiamo a fondo la tecnologia dell'“espansione memoria”.}

\booksubsection{Definizione essenziale}

L'“espansione memoria” dei cellulari Android (chiamata anche estensione memoria, fusione memoria, RAM Plus, ecc.)\textbf{non aggiunge fisicamente chip di memoria RAM}, ma è una tecnologia software a livello di sistema: converte dinamicamente parte della memoria interna (ROM/flash UFS) in memoria RAM virtuale.

Per fare un paragone: la RAM fisica è il “banco di lavoro ad alta velocità”, l'espansione memoria è il “magazzino temporaneo”. Quando il banco non basta, il sistema sposta i dati temporaneamente inutilizzati nel magazzino e li riprende quando servono. Il magazzino ha grande capacità, ma gli spostamenti sono sicuramente più lenti di averli direttamente sul banco.

\booksubsection{Principio tecnico}

I principali produttori usano due approcci:

\subsubsection*{Compressione zRAM}

All'interno della RAM viene ricavato uno spazio dati compresso in cui conservare compresso il codice delle app inattive, con un ratio di compressione di circa 2:1. I dati restano in RAM, veloci e senza usura della memoria flash, ma la CPU deve partecipare alla compressione/decompressione.

\subsubsection*{Partizione swap}

Si divide direttamente uno spazio fisso nella memoria flash UFS da usare come memoria virtuale, scrivendo i dati dei processi in background nello storage. Dipende dalla larghezza di banda dello storage, con velocità di lettura/scrittura di gran lunga inferiori alla memoria fisica.

Il “12GB+8GB RAM” reclamizzato dai produttori è in genere la combinazione di 12GB di RAM fisica + 8GB di memoria virtuale su flash. Vale la pena notare che zRAM di default non occupa storage del telefono: solo quando zRAM è completamente piena di dati il sistema crea un file di swap sullo storage, situazione che si verifica raramente.

\booksubsection{Il divario di velocità è il collo di bottiglia chiave}

La limitazione fondamentale di questa tecnologia è il divario di velocità. L'attuale memoria flash UFS 4.0 mainstream ha una velocità di lettura sequenziale di circa 4200 MB/s, solo 1/10-1/15 della memoria fisica LPDDR5X. Pertanto, nella fase di compressione zRAM a basso carico, la velocità è vicina a quella della RAM reale; ma in scenari ad alto carico, la velocità crolla a poche decine di volte inferiore.

\booksubsection{Scenari applicativi e confini dei benefici}

L'espansione memoria non è onnipotente; l'effetto varia in base al modello e allo scenario d'uso:

\begin{itemize}
  \item \textbf{Casi utili}: modelli con 6GB di RAM fisica + UFS 3.1 o superiore, spazio residuo $\geq$30GB; attivando 3GB di espansione, le app in background sopravvivono in 3-4 in più, il tasso di reload si riduce del 68\%, adatto a utenti multi-tasking intensivo come rider delle consegne o autisti di ride-hailing.
  \item \textbf{Scenari quotidiani}: miglioramenti significativi per multi-account WeChat, mantenimento delle schede del browser, risposta ai messaggi istantanei; il tempo di avvio a freddo delle app in background si riduce notevolmente.
  \item \textbf{Scenari limitati}: miglioramenti limitati per il caricamento di giochi pesanti, il montaggio video 4K e altri scenari ad alta larghezza di banda.
  \item \textbf{Casi dannosi}: modelli con 8GB di RAM + UFS 2.2 e spazio residuo insufficiente; dopo l'attivazione, la fluctuazione del frame rate nei giochi aumenta, lo scorrimento dei video brevi consuma più batteria e il cambio app scatta.
  \item \textbf{Casi inefficaci}: modelli con 16GB di RAM o più beneficiano quasi zero, al massimo una app in più in background, ma aumentano il carico CPU del 12\% e il calore di 2-3°C. Per modelli con RAM superiore a 12GB, si consiglia di disattivarla per evitare calcoli di compressione inutili.
\end{itemize}

\booksubsection{Denominazione dei produttori e modalità di attivazione}

I diversi brand la chiamano in modi diversi: Xiaomi “Espansione memoria”, Huawei “Fusione memoria”, Samsung “RAM Plus”, vivo “Extended RAM”, OnePlus “Memory Extension”. È una capacità nativa del sistema, non legata direttamente alla versione sottostante di Android.

Il percorso tipico di attivazione: entra in [Impostazioni] [Altre impostazioni] trova la voce [Memoria] o [Prestazioni] attiva “Espansione memoria” e scegli la capacità (in genere 2GB-7GB), riavvia per rendere effettiva. Dopo l'attivazione, il sistema mostrerà una combinazione di memoria come “8GB+4GB”.

\begin{tcbwarn}{Avvertenze}
\begin{itemize}
  \item L'espansione memoria occupa parte dello storage del telefono; assicurati che lo storage residuo sia almeno 20-30GB, altrimenti potrebbe compromettere la fluidità del sistema.
  \item L'uso della memoria flash come memoria virtuale comporta usura da lettura/scrittura; anche se la memoria flash UFS moderna ha vita lunga, swap ad alta frequenza a lungo termine può accelerare l'invecchiamento dello storage.
  \item Se il cellulare ha poca RAM fisica (4GB o meno), l'espansione memoria è solo un palliativo e non risolve alla radice il problema dei lag.
\end{itemize}

\end{tcbwarn}

\begin{tcbinfo}{In sintesi}
L'espansione memoria è una tecnologia matura di ottimizzazione a livello sistema, che può migliorare le capacità multi-tasking e la sopravvivenza in background dei modelli medio-bassi senza modificare l'hardware, ma è limitata dal divario di velocità tra flash e RAM fisica e non è una “pozione magica”. Se attivarla o meno e con quale capacità, dipende dalla dimensione della RAM fisica del modello, dalle specifiche della flash, dallo spazio residuo e dalle abitudini d'uso.

\end{tcbinfo}

\begin{tcbpractice}{Fai pratica: verifica se il tuo cellulare ha l'espansione memoria}
\begin{enumerate}
  \item Apri [Impostazioni], cerca "Espansione memoria", "Fusione memoria", "RAM Plus" o "Extended RAM".
  \item Una volta trovata, verifica se è attiva e quale capacità è impostata.
  \item Confronta con i contenuti di questa sezione: quanta RAM fisica ha il tuo modello? Quali specifiche flash? Lo spazio residuo è sufficiente?
  \item Prova a disattivarla o a ridurla, osserva per una settimana i cambiamenti nella sopravvivenza in background e nel riscaldamento, trova l'impostazione più adatta a te.
\end{enumerate}

\end{tcbpractice}
